\chapter{Rasa Suka dan Kagum: Fondness \& Admiration}

\begin{insightbox}
\textit{``Fondness and admiration adalah \highlight{antidot paling kuat} untuk contempt. Ketika kamu secara aktif mencari dan mengungkapkan hal-hal yang kamu hargai dari pasanganmu, kamu membangun fondasi yang tidak bisa digoyangkan oleh konflik sekalipun.''}
\end{insightbox}

\section{Memahami Sistem Fondness dan Admiration}

\subsection{Apa itu Fondness System?}

\keyterm{Fondness and admiration system} adalah fondasi psikologis yang mendasari setiap hubungan yang sehat. Sistem ini terdiri dari \highlight{perasaan hangat}, \highlight{rasa hormat}, dan \highlight{kagum} yang kamu miliki terhadap pasanganmu.

Dr. John Gottman menemukan bahwa pasangan yang mempertahankan sistem fondness yang kuat memiliki kemungkinan \textbf{86\% lebih tinggi} untuk tetap bersama dalam jangka panjang.

\begin{praktikbox}[Definisi Fondness System]
Sistem fondness adalah:
\begin{itemize}
    \item \textbf{Memori positif} yang kamu simpan tentang pasangan
    \item \textbf{Rasa syukur} atas keberadaannya dalam hidupmu
    \item \textbf{Kagum} pada kualitas dan kekuatan karakternya
    \item \textbf{Kebanggaan} menjadi pasangannya
    \item \textbf{Kasih sayang} yang tulus dan tidak bersyarat
\end{itemize}
\end{praktikbox}

\subsection{Fondness sebagai Antidot untuk Contempt}

Contempt (penghinaan) adalah \highlight{Death Horseman} paling berbahaya dalam hubungan. Tapi fondness adalah \textbf{vaksin} yang mencegah contempt muncul.

\begin{center}
\begin{tikzpicture}[
    box/.style={draw=#1, line width=2pt, rounded corners=15pt,
                fill=#1!10, minimum width=5cm, minimum height=3cm, align=center}
]
    % Fondness Box
    \node[box=sagegreen] (fond) at (0,0) {
        \textbf{\large FONDNESS}\\[0.3cm]
        \small
        {\small
        \begin{tabular}{ll}
        \textcolor{sagegreen}{$\bullet$} & Apresiasi \\
        \textcolor{sagegreen}{$\bullet$} & Kagum \\
        \textcolor{sagegreen}{$\bullet$} & Rasa syukur \\
        \textcolor{sagegreen}{$\bullet$} & Kenangan indah \\
        \end{tabular}
        }
    };
    
    % Shield
    \node[draw=primary, line width=3pt, rounded corners=10pt,
          fill=primary!10, minimum width=3cm, minimum height=1.5cm] (shield) at (5,0) {
        \textbf{\large PERISAI}
    };
    
    % Contempt Box
    \node[box=accent] (cont) at (10,0) {
        \textbf{\large CONTEMPT}\\[0.3cm]
        \small
        {\small
        \begin{tabular}{ll}
        \textcolor{accent}{$\bullet$} & Menghina \\
        \textcolor{accent}{$\bullet$} & Meremehkan \\
        \textcolor{accent}{$\bullet$} & Sarkasme \\
        \textcolor{accent}{$\bullet$} & Superioritas \\
        \end{tabular}
        }
    };
    
    % Arrows
    \draw[->, line width=2pt, sagegreen] (fond) -- (shield);
    \draw[->, line width=2pt, accent] (cont) -- (shield);
    \draw[->, line width=3pt, sagegreen] (shield.east) -- ++(1,0);
\end{tikzpicture}
\end{center}

\section{Ilmu di Balik Fondness: Prediktor Stabilitas Hubungan}

\subsection{Penelitian Gottman yang Mengubah Segalanya}

Dalam penelitian longitudinal selama \textbf{16 tahun}, Gottman mengamati ribuan pasangan dan menemukan pola yang menakjubkan:

\begin{center}
\renewcommand{\arraystretch}{1.5}
\begin{tabular}{|>{\raggedright}p{4cm}|>{\raggedright}p{5cm}|>{\raggedright\arraybackslash}p{5cm}|}
\hline
\rowcolor{primary!20}
\textbf{Indikator} & \textbf{Pasangan Bahagia} & \textbf{Pasangan Tidak Bahagia} \\
\hline
Ekspresi fondness/hari & 5-20 kali & 0-2 kali \\
\hline
Rasio positif-negatif & 20:1 (konflik) & 0,8:1 \\
\hline
Kenangan masa awal & Hangat, detail & Dingin, kabur \\
\hline
Respons terhadap keberhasilan pasangan & Active Constructive & Passive/Destructive \\
\hline
\end{tabular}
\end{center}

\begin{insightbox}
\textbf{Fakta Mengejutkan:} Pasangan yang tetap bahagia selama 20+ tahun tidak memiliki \textit{lebih sedikit} konflik. Mereka hanya memiliki \highlight{lebih banyak fondness} yang melindungi hubungan saat konflik terjadi.
\end{insightbox}

\subsection{Neuroscience of Fondness}

Ketika kamu merasakan fondness:

\begin{enumerate}
    \item \textbf{Oxytocin naik} --- memperkuat ikatan emosional
    \item \textbf{Cortisol turun} --- mengurangi stres
    \item \textbf{Prefrontal cortex aktif} --- kemampuan regulasi emosi meningkat
    \item \textbf{Amygdala tenang} --- deteksi ``ancaman'' berkurang
\end{enumerate}

Ini menciptakan \keyterm{upward spiral} --- semakin kamu merasakan fondness, semakin mudah untuk melihat kebaikan pasangan.

\section{Mengapa Fondness Memudar?}

\subsection{Erosi Normal vs Penurunan Berbahaya}

\begin{warningbox}[Bedakan Dua Jenis Penurunan]
\textbf{Erosi Normal (Natural Decline):}
\begin{itemize}
    \item Terjadi perlahan selama 2-3 tahun pertama
    \item Fondness tetap ada, hanya frekuensi ungkapan berkurang
    \item Masih mudah mengingat hal-hal positif saat diminta
    \item Tidak ada rasa jengkel mendalam
\end{itemize}

\textbf{Penurunan Berbahaya (Dangerous Decline):}
\begin{itemize}
    \item Fondness digantikan oleh kritik berulang
    \item Sulit mengingat hal positif tentang pasangan
    \item Rasa jengkel muncul bahkan saat tidak ada masalah
    \item Kenangan masa lalu diceritakan dengan negatif
\end{itemize}
\end{warningbox}

\subsection{Faktor yang Mempercepat Hilangnya Fondness}

\begin{center}
\begin{tikzpicture}[
    factor/.style={draw=accent, line width=1.5pt, rounded corners=8pt,
                   fill=accent!10, minimum width=4cm, minimum height=1.8cm,
                   align=center, font=\small}
]
    % Title
    \node[font=\Large\bfseries, primary] at (0,4) {Faktor Percepatan Erosi Fondness};
    
    % Factors
    \node[factor] (stress) at (-4,2) {\textbf{Stres Kronis}\\Kerja, finansial, kesehatan};
    \node[factor] (neglect) at (0,2) {\textbf{Neglect}\\Tidak ada quality time};
    \node[factor] (conflict) at (4,2) {\textbf{Konflik Tak terselesaikan}\\Resentment menumpuk};
    \node[factor] (comparison) at (-4,0) {\textbf{Perbandingan}\\Media sosial, ekspektasi tidak realistis};
    \node[factor] (taking) at (0,0) {\textbf{Menganggap Remeh}\\Merasa ``sudah pasti''};
    \node[factor] (external) at (4,0) {\textbf{Fokus Eksternal}\\Anak, pekerjaan, orang tua};
    
    % Arrow down
    \node[draw=primary, line width=3pt, rounded corners=10pt,
          fill=primary!20, minimum width=10cm, minimum height=1cm] (result) at (0,-2) {
        \textbf{\large FONDNESS MENURUN DRASTIS}
    };
    
    % Connecting lines
    \draw[->, line width=1pt] (stress) -- (result);
    \draw[->, line width=1pt] (neglect) -- (result);
    \draw[->, line width=1pt] (conflict) -- (result);
    \draw[->, line width=1pt] (comparison) -- (result);
    \draw[->, line width=1pt] (taking) -- (result);
    \draw[->, line width=1pt] (external) -- (result);
\end{tikzpicture}
\end{center}

\section{Self-Assessment: Mengukur Level Fondness-mu}

\begin{exercisebox}[Fondness Assessment Scale]
\textbf{Petunjuk:} Jawab setiap pertanyaan dengan skala 1-5:
\begin{itemize}
    \item[1] = Sangat tidak setuju
    \item[3] = Netral
    \item[5] = Sangat setuju
\end{itemize}

\subsection*{Bagian A: Ungkapan Fondness}

\begin{enumerate}
    \item Saya sering memuji pasangan saya secara spontan.
    \item Saya merasa bangga ketika menceritakan pasangan saya kepada orang lain.
    \item Saya menunjukkan kasih sayang fisik setiap hari (peluk, cium, gandeng tangan).
    \item Saya mengucapkan ``terima kasih'' atas hal-hal kecil yang pasangan lakukan.
    \item Saya sering tersenyum saat melihat/membayangkan pasangan saya.
\end{enumerate}

\subsection*{Bagian B: Pikiran dan Perasaan}

\begin{enumerate}
    \setcounter{enumi}{5}
    \item Ketika terpisah, saya sering memikirkan hal positif tentang pasangan.
    \item Saya merasa beruntung memiliki pasangan saya.
    \item Saya bisa dengan mudah menyebutkan 5 kualitas yang saya kagumi dari pasangan.
    \item Kenangan awal hubungan kami masih membawa perasaan hangat.
    \item Saya percaya pasangan saya adalah orang yang baik, meski punya kekurangan.
\end{enumerate}

\subsection*{Bagian C: Saat Konflik}

\begin{enumerate}
    \setcounter{enumi}{10}
    \item Saat marah, saya masih bisa mengingat hal-hal baik tentang pasangan.
    \item Saya menghindari menghina atau meremehkan pasangan saat bertengkar.
    \item Setelah konflik, saya bisa kembali merasa hangat pada pasangan.
    \item Saya tidak membuat kesalahan kecil menjadi masalah besar.
    \item Saya bisa memaafkan dan melupakan dengan relatif mudah.
\end{enumerate}

\subsection*{Scoring}

\begin{center}
\renewcommand{\arraystretch}{1.4}
\begin{tabular}{|c|c|l|}
\hline
\rowcolor{primary!20}
\textbf{Skor Total} & \textbf{Kategori} & \textbf{Interpretasi} \\
\hline
60-75 & \textcolor{sagegreen}{\textbf{Sangat Tinggi}} & Fondness system sangat kuat. Pertahankan! \\
\hline
45-59 & \textcolor{primary}{\textbf{Tinggi}} & Fondness baik, tapi ada ruang untuk ditingkatkan. \\
\hline
30-44 & \textcolor{accent}{\textbf{Sedang}} & Fondness mulai memudar. Perlu perhatian serius. \\
\hline
15-29 & \textcolor{red}{\textbf{Rendah}} & \textbf{Warning!} Fondness berbahaya. Butuh intervensi segera. \\
\hline
\end{tabular}
\end{center}
\end{exercisebox}

\section{Membangun Kembali Fondness: Proses Langkah demi Langkah}

\subsection{Untuk Pasangan yang Masih Memiliki Fondness}

\begin{tipbox}[Strategi Maintenance]
\begin{enumerate}
    \item \textbf{Buat Ritual Ungkapan}
    \begin{itemize}
        \item Setiap pagi: sebutkan 1 hal yang kamu hargai
        \item Setiap malam: ucapkan terima kasih untuk 1 hal spesifik hari ini
    \end{itemize}
    \item \textbf{Ciptakan Tradisi Positif}
    \begin{itemize}
        \item Anniversary celebration yang bermakna
        \item Weekly date night tanpa gadget
        \item Monthly reflection bersama
    \end{itemize}
    \item \textbf{Dokumentasikan Momen Bahagia}
    \begin{itemize}
        \item Foto album digital
        \item Journal kenangan bersama
        \item Box memorabilia
    \end{itemize}
\end{enumerate}
\end{tipbox}

\subsection{Untuk Pasangan yang Kehilangan Fondness}

\begin{warningbox}[Proses Rebuilding Butuh Waktu]
Membangun kembali fondness bukanlah proses instan. Dibutuhkan \textbf{6-12 bulan} konsistensi untuk mengembalikan sistem fondness yang kuat.
\end{warningbox}

\textbf{Langkah 1: Akui Masalahnya}
\begin{itemize}
    \item ``Aku sadar aku sudah lama tidak menunjukkan apresiasi padamu. Aku ingin mengubahnya.''
    \item Hindari defensive atau menyalahkan
    \item Terima bahwa ini butuh usaha dari kedua belah pihak
\end{itemize}

\textbf{Langkah 2: Mulai dengan Mengamati}
\begin{itemize}
    \item Luangkan 5 menit setiap hari untuk ``meneliti'' pasanganmu
    \item Catat hal-hal kecil yang dia lakukan dengan baik
    \item Perhatikan kualitas karakternya yang mungkin terlupakan
\end{itemize}

\textbf{Langkah 3: Ungkapkan Secara Bertahap}
\begin{itemize}
    \item Mulai dengan yang aman dan umum
    \item Tingkatkan spesifisitas dan frekuensi perlahan
    \item Jangan berharap respons positif langsung --- berikan waktu
\end{itemize}

\section{Praktik Harian: 20+ Tindakan Spesifik}

\subsection{Verbal Appreciations (Ungkapan Verbal)}

\begin{praktikbox}[Contoh Ungkapan Verbal]
\textbf{Tingkat 1 - Umum:}
\begin{itemize}
    \item ``Terima kasih ya.''
    \item ``Aku appreciate yang kamu lakukan.''
    \item ``Kamu hebat.''
\end{itemize}

\textbf{Tingkat 2 - Spesifik:}
\begin{itemize}
    \item ``Terima kasih sudah membuatkan kopi pagi ini. Aku butuh itu.''
    \item ``Aku notice kamu sudah bekerja keras minggu ini. Aku bangga padamu.''
    \item ``Cara kamu menangani situasi tadi sangat bijaksana.''
\end{itemize}

\textbf{Tingkat 3 - Mendalam:}
\begin{itemize}
    \item ``Aku bersyukur setiap hari karena kamu adalah pasanganku.''
    \item ``Kualitas kepedulianmu adalah hal yang paling aku kagumi dari dirimu.''
    \item ``Aku jadi orang yang lebih baik karena kamu ada di hidupku.''
\end{itemize}
\end{praktikbox}

\subsection{Affection Fisik}

\begin{center}
\renewcommand{\arraystretch}{1.4}
\begin{tabular}{|>{\raggedright}p{3cm}|>{\raggedright}p{4.5cm}|>{\raggedright\arraybackslash}p{6cm}|}
\hline
\rowcolor{sagegreen!20}
\textbf{Jenis Sentuhan} & \textbf{Contoh} & \textbf{Kapan Menggunakannya} \\
\hline
Micro-affection & Tap punggung, squeeze tangan, touch shoulder & Saat lewat, saat bekerja sama, saat berpisah \\
\hline
Greeting affection & Peluk saat pulang, cium dahi & Setiap kali bertemu setelah terpisah \\
\hline
Comforting touch & Peluk lama, gosok punggung & Saat pasangan stres, sedih, atau capek \\
\hline
Playful touch & Bercanda, bergandengan jalan & Saat santai, jalan-jalan, momen ringan \\
\hline
Intimate touch & Ciuman, berpelukan tanpa agenda & Saat quality time, sebelum tidur \\
\hline
\end{tabular}
\end{center}

\subsection{Acts of Service (Tindakan Pelayanan)}

\begin{enumerate}
    \item \textbf{Morning routine support:} Siapkan sarapan atau kopi
    \item \textbf{Evening wind-down:} Urut bahu pasangan setelah hari yang panjang
    \item \textbf{Task takeover:} Ambil alih tugas yang biasanya pasangan lakukan
    \item \textbf{Proactive help:} Lihat apa yang butuh dikerjakan dan lakukan tanpa diminta
    \item \textbf{Comfort creation:} Siapkan mandi hangat, nyalakan lilin, atmosfer nyaman
\end{enumerate}

\subsection{Quality Time}

\begin{tipbox}[Ide Quality Time 15-30 Menit]
\begin{itemize}
    \item[$\bullet$] \textbf{Morning check-in:} Duduk bersama sambil minum kopi, bicara rencana hari
    \item[$\bullet$] \textbf{Evening debrief:} Ceritakan satu hal baik dan satu tantangan dari hari ini
    \item[$\bullet$] \textbf{Walk together:} Jalan sore sambil bergandengan
    \item[$\bullet$] \textbf{Cook together:} Masak bersama sambil bercanda
    \item[$\bullet$] \textbf{Tech-free time:} Duduk berhadapan tanpa gadget selama 20 menit
\end{itemize}
\end{tipbox}

\subsection{Small Gestures (Hadiah Kecil)}

\begin{enumerate}
    \item Catatan singkat di tempat tak terduga (kotak bekal, cermin kamar mandi)
    \item Membeli makanan favoritnya saat pulang kerja
    \item Kirim pesan singkat di tengah hari: ``Lagi mikirin kamu."
    \item Ambilkan benda yang dia butuhkan tanpa diminta
    \item Buat playlist lagu yang mengingatkan padanya
    \item Foto momen kecil dan kirim dengan caption manis
    \item Siapkan ``emergency kit'' di mobil: makanan favorit, obat, tissue
\end{enumerate}

\section{The 7-Week Fondness Exercise: Rencana Mingguan}

\begin{exercisebox}[Program 7 Minggu Transformasi]
\renewcommand{\arraystretch}{1.6}
\begin{tabular}{|c|>{\raggedright}p{4cm}|>{\raggedright\arraybackslash}p{7cm}|}
\hline
\rowcolor{primary!20}
\textbf{Minggu} & \textbf{Fokus} & \textbf{Tugas Harian} \\
\hline
1 & \textbf{Apresiasi Dasar} & Ucapkan 1 hal spesifik yang kamu hargai setiap hari \\
\hline
2 & \textbf{Kenangan Indah} & Ceritakan ulang 1 kenangan indah masa pacaran/dating \\
\hline
3 & \textbf{Kualitas Kagum} & Identifikasi dan ucapkan 1 kualitas karakter yang kamu kagumi \\
\hline
4 & \textbf{Mikro-Momen} & Ciptakan 3 micro-momen koneksi (sentuhan, eye contact, senyum) \\
\hline
5 & \textbf{Public Pride} & Ungkapkan kebanggaan pada pasangan di depan orang lain 1x \\
\hline
6 & \textbf{Love Notes} & Tulis catatan singkat/manis yang ditinggalkan di tempat tak terduga \\
\hline
7 & \textbf{Integrasi} & Lakukan semua praktik minggu 1-6 setiap hari \\
\hline
\end{tabular}

\vspace{0.5cm}
\textbf{Tips Keberhasilan:}
\begin{itemize}
    \item Jadwalkan waktu yang sama setiap hari (misal: saat sarapan dan sebelum tidur)
    \item Gunakan reminder di HP untuk 30 hari pertama
    \item Buat journal tracking untuk mencatat progress
    \item Rayakan milestone bersama setiap minggu
\end{itemize}
\end{exercisebox}

\section{Mengungkapkan Admiration dengan Efektif}

\subsection{Specific vs Generic Compliments}

\begin{warningbox}[Hindari Pujian Generik]
Pujian seperti ``Kamu baik'' atau ``Kamu cantik/ganteng'' memiliki dampak terbatas karena:
\begin{itemize}
    \item Terlalu umum --- bisa diucapkan ke siapa saja
    \item Tidak menunjukkan bahwa kamu benar-benar \textit{melihat} pasanganmu
    \item Mudah dilupakan
\end{itemize}
\end{warningbox}

\begin{praktikbox}[Formula Pujian Spesifik]
\textbf{Template:} ``Aku [perasaan] ketika kamu [tindakan spesifik] karena [dampak/makna].''

\textbf{Contoh:}
\begin{itemize}
    \item ``Aku \textbf{merasa dihargai} ketika kamu \textbf{menanyakan hari ku hari ini} karena \textbf{itu menunjukkan kamu peduli dengan apa yang kulakukan}.''
    \item ``Aku \textbf{bangga} melihat cara kamu \textbf{menangani deadline yang ketat minggu ini} karena \textbf{kamu tetap tenang dan profesional}.''
    \item ``Aku \textbf{terkesan} dengan \textbf{kesabaranmu saat mengajari anak belajar} karena \textbf{kamu menjelaskan dengan lembut meski capek}.''
\end{itemize}
\end{praktikbox}

\subsection{The 5:1 Ratio dalam Action}

Gottman menemukan bahwa pasangan bahagia memiliki rasio \textbf{5 interaksi positif : 1 interaksi negatif} saat konflik.

\begin{center}
\begin{tikzpicture}[
    posbox/.style={draw=sagegreen, line width=2pt, rounded corners=5pt,
                fill=sagegreen!20, minimum width=1.5cm, minimum height=1cm, font=\small},
    negbox/.style={draw=accent, line width=2pt, rounded corners=5pt,
                fill=accent!20, minimum width=1.5cm, minimum height=1cm, font=\small}
]
    % 5 positive boxes
    \node[posbox] (p1) at (0,0) {+};
    \node[posbox] (p2) at (2,0) {+};
    \node[posbox] (p3) at (4,0) {+};
    \node[posbox] (p4) at (6,0) {+};
    \node[posbox] (p5) at (8,0) {+};
    \node[negbox] (n1) at (10,0) {-};
    
    % Label
    \node[font=\small, sagegreen] at (4,1) {5 Interaksi Positif};
    \node[font=\small, accent] at (10,1) {1 Interaksi Negatif};
    
    % Ratio
    \node[draw=primary, line width=3pt, rounded corners=10pt,
          fill=primary!10, minimum width=4cm, minimum height=1cm] at (5,-2) {
        \textbf{Rasio 5:1 = Hubungan Sehat}
    };
\end{tikzpicture}
\end{center}

\section{Menerima Admiration dengan Grace}

\subsection{Kenapa Sulit Menerima Pujian?}

Banyak orang merasa \textbf{tidak nyaman} saat dipuji karena:
\begin{itemize}
    \item \textbf{Low self-esteem:} Merasa tidak pantas dipuji
    \item \textbf{Fear of expectations:} Takut harus mempertahankan standar
    \item \textbf{Cultural conditioning:} Dibiasakan untuk merendah
    \item \textbf{Distrust:} Meragukan kejujuran pujian
\end{itemize}

\subsection{Cara Menerima dengan Baik}

\begin{tipbox}[Respons yang Membangun]
\textbf{Lakukan:}
\begin{enumerate}
    \item \textbf{Terima dengan sederhana:} ``Terima kasih, aku appreciate pujianmu.''
    \item \textbf{Akui usaha:} ``Aku memang berusaha keras untuk itu.''
    \item \textbf{Bagikan credit:} ``Aku belajar banyak dari kamu juga.''
    \item \textbf{Reflect back:} ``Itu berarti banyak dari kamu.''
\end{enumerate}

\textbf{Hindari:}
\begin{itemize}
    \item Menyangkal: ``Ah, enggak kok, biasa aja.''
    \item Defleksi: ``Bajuku aja yang bagus.''
    \item Counter-criticism: ``Tapi aku jelek di hal lain.''
    \item Silence awkward: Diam saja tanpa respons
\end{itemize}
\end{tipbox}

\section{Menjaga Fondness Saat Konflik}

\subsection{Softened Start-up dengan Fondness}

Saat membahas masalah sulit, mulailah dengan fondness:

\begin{center}
\renewcommand{\arraystretch}{1.5}
\begin{tabular}{|>{\raggedright}p{6cm}|>{\raggedright\arraybackslash}p{8cm}|}
\hline
\rowcolor{accent!20}
\textbf{Hard Start-up (Negatif)} & \textbf{Softened Start-up (Dengan Fondness)} \\
\hline
``Kamu selalu lupa membuang sampah!'' & ``Aku tahu kamu sibuk, tapi aku butuh bantuan dengan sampah.'' \\
\hline
``Kamu nggak pernah dengerin aku!'' & ``Aku tahu kamu care, tapi saat ini aku butuh full attention-mu.'' \\
\hline
``Kamu egois banget sih!'' & ``Biasanya kamu very considerate, tapi keputusan ini bikin aku kecewa.'' \\
\hline
\end{tabular}
\end{center}

\subsection{Repair Attempts yang Efektif}

\begin{praktikbox}[Repair Attempts dengan Fondness]
\textbf{Phrase untuk mengembalikan fondness saat konflik:}
\begin{itemize}
    \item ``Aku tahu kita bisa lewati ini. Aku di sini.''
    \item ``Ini bukan tentang kita melawan satu sama lain, tapi kita vs masalah ini.''
    \item ``Aku marah, tapi aku tetap cinta kamu.''
    \item ``Aku appreciate kamu masih mau bicara tentang ini.''
    \item ``Maaf, aku salah. Kamu pantas mendapatkan pendengaran yang lebih baik.''
\end{itemize}
\end{praktikbox}

\section{Integrasi dengan Attachment Styles}

\begin{insightbox}
Setiap attachment style memiliki cara unik dalam memberikan dan menerima fondness. Memahami ini membantu pasangan saling memenuhi kebutuhan emosional dengan lebih efektif.
\end{insightbox}

\subsection{Tabel Integrasi Lengkap}

\begin{center}
\renewcommand{\arraystretch}{1.6}
\begin{tabular}{|>{\raggedright}p{2.5cm}|>{\raggedright}p{4cm}|>{\raggedright}p{4cm}|>{\raggedright\arraybackslash}p{4cm}|}
\hline
\rowcolor{primary!20}
\textbf{Attachment} & \textbf{Cara Memberi Fondness} & \textbf{Cara Menerima Fondness} & \textbf{Kebutuhan Spesifik} \\
\hline
\textbf{Anchor} & Natural dan berimbang; verbal, fisik, service dengan mudah & Nyaman dan menerima dengan terbuka & Konsistensi dan kualitas waktu \\
\hline
\textbf{Island} & Melalui acts of service lebih dari verbal; menunjukkan dengan tindakan & Sulit menerima verbal, lebih nyaman dengan non-verbal & Jangan terlalu ``tebal'', beri ruang \\
\hline
\textbf{Wave} & Generous dan banyak; verbal dan fisik melimpah & Butuh validasi berulang, takut pujian palsu & Reassurance dan konsistensi \\
\hline
\end{tabular}
\end{center}

\subsection{Strategi untuk Setiap Kombinasi}

\begin{tipbox}[Panduan Berdasarkan Kombinasi]
\textbf{Anchor + Island:}
\begin{itemize}
    \item Anchor: Hargai cara non-verbal Island menunjukkan cinta
    \item Island: Latih diri untuk menerima dan merespons verbal affection
\end{itemize}

\textbf{Anchor + Wave:}
\begin{itemize}
    \item Anchor: Berikan reassurance ekstra tanpa merasa ``cukup sudah''
    \item Wave: Terima bahwa Anchor tidak perlu ``diuji'' terus-menerus
\end{itemize}

\textbf{Island + Wave:}
\begin{itemize}
    \item Island: Berusaha lebih ekspresif meski tidak nyaman
    \item Wave: Beri ruang saat Island perlu withdraw, jangan personalisasi
\end{itemize}
\end{tipbox}

\section{Studi Kasus: Rebuilding dari Nol}

\subsection{Kasus 1: David dan Sarah (15 Tahun Menikah)}

\begin{insightbox}
\textbf{Situasi:} David dan Sarah hampir bercerai. David merasa Sarah ``tidak pernah puas'', Sarah merasa David ``tidak peduli''. Fondness hampir nol --- mereka tidak pernah memuji satu sama lain selama 3 tahun terakhir.

\textbf{Intervensi:}
\begin{enumerate}
    \item \textbf{Minggu 1-2:} Mereka hanya diminta untuk \textit{mencatat} 1 hal positif tentang pasangan setiap hari, tanpa mengungkapkan.
    \item \textbf{Minggu 3-4:} Mulai mengungkapkan 1 apresiasi sederhana setiap hari.
    \item \textbf{Minggu 5-6:} Share kenangan indah dari 15 tahun pernikahan.
    \item \textbf{Minggu 7-8:} Weekly date night dengan tema ``apresiasi''.
\end{enumerate}

\textbf{Hasil:} Setelah 6 bulan, mereka melaporkan ``jatuh cinta lagi''. David belajar melihat usaha Sarah yang dia anggap ``mengeluh'' sebenarnya adalah kepedulian. Sarah belajar melihat ketenangan David bukan ``tidak peduli'' tapi ``stabilitas''.
\end{insightbox}

\subsection{Kasus 2: Mike dan Lisa (Island + Wave Dynamic)}

\begin{insightbox}
\textbf{Situasi:} Mike (Island) dan Lisa (Wave) dalam spiral negatif. Lisa mendesak affection, Mike semakin menjauh. Lisa merasa tidak dicintai, Mike merasa tidak cukup baik.

\textbf{Strategi Khusus:}
\begin{itemize}
    \item Mike diajari ``micro-affection'' --- sentuhan kecil yang tidak membuatnya overwhelm
    \item Lisa diajari ``receive without pursuit'' --- menerima affection tanpa langsung menuntut lebih
    \item Mereka membuat ``affection schedule'' yang bisa diprediksi
\end{itemize}

\textbf{Breakthrough:} Ketika Lisa berhenti mengejar, Mike justru mulai mendekat. Ketika Mike memberi affection tanpa diminta, Lisa merasa aman.
\end{insightbox}

\section{30-Day Fondness Challenge}

\begin{exercisebox}[Tantangan 30 Hari: Prompt Harian]
\textbf{Instruksi:} Setiap hari, gunakan prompt berikut untuk mengungkapkan fondness pada pasangan. Variasikan cara ungkapannya (verbal, catatan, tindakan).

\renewcommand{\arraystretch}{1.4}
\begin{tabular}{|c|>{\raggedright\arraybackslash}p{10cm}|}
\hline
\rowcolor{primary!20}
\textbf{Hari} & \textbf{Prompt} \\
\hline
1 & Sebutkan 1 hal yang pasanganmu lakukan hari ini yang kamu hargai. \\
\hline
2 & Ceritakan kenangan pertama kalian dengan detail. \\
\hline
3 & Apa 1 kualitas karakter pasanganmu yang kamu kagumi? \\
\hline
4 & Ucapkan terima kasih untuk hal kecil yang sering dianggap remeh. \\
\hline
5 & Peluk pasanganmu selama 20 detik tanpa bicara. \\
\hline
6 & Tulis catatan singkat dan sembunyikan di tempat tak terduga. \\
\hline
7 & Ceritakan pada orang lain (di depan pasanganmu) hal baik tentangnya. \\
\hline
8 & Lakukan 1 tugas yang biasanya pasanganmu kerjakan. \\
\hline
9 & Sebutkan 3 hal yang membuatmu tertawa bersamanya. \\
\hline
10 & Buat playlist lagu yang mengingatkan padanya. \\
\hline
11 & Berikan komplimen spesifik tentang penampilannya hari ini. \\
\hline
12 & Ceritakan momen saat kamu paling bangga padanya. \\
\hline
13 & Kirim pesan singkat di tengah hari: ``Sedang mikirin kamu.'' \\
\hline
14 & Lakukan aktivitas favoritnya meski bukan favoritmu. \\
\hline
15 & Sebutkan 1 hal yang kamu pelajari darinya. \\
\hline
\end{tabular}
\end{exercisebox}

\begin{exercisebox}[Tantangan 30 Hari: Lanjutan]
\renewcommand{\arraystretch}{1.4}
\begin{tabular}{|c|>{\raggedright\arraybackslash}p{10cm}|}
\hline
\rowcolor{primary!20}
\textbf{Hari} & \textbf{Prompt} \\
\hline
16 & Buat daftar 10 hal yang kamu syukuri tentang pernikahan/hubunganmu. \\
\hline
17 & Berikan massage sederhana di bahu atau punggung. \\
\hline
18 & Ceritakan kenangan lucu yang melibatkan kalian berdua. \\
\hline
19 & Ucapkan ``Aku cinta kamu'' dengan cara yang berbeda dari biasanya. \\
\hline
20 & Siapkan makanan favoritnya sebagai kejutan. \\
\hline
21 & Sebutkan 1 mimpi/karir goal-nya dan dukung secara verbal. \\
\hline
22 & Tulis surat singkat (setengah halaman) tentang kenapa kamu bersyukur. \\
\hline
23 & Lakukan eye contact selama 2 menit tanpa bicara. \\
\hline
24 & Berikan hadiah kecil tanpa alasan khusus. \\
\hline
25 & Ceritakan bagaimana pasanganmu telah membantumu tumbuh. \\
\hline
26 & Ucapkan maaf untuk hal kecil yang mungkin menyakitinya. \\
\hline
27 & Buat ritual selamat datang saat dia pulang dengan affection. \\
\hline
28 & Sebutkan 1 hal yang kamu nantikan tentang masa depan bersama. \\
\hline
29 & Lakukan ``debriefing positif'': ceritakan 1 hal baik dari minggu ini. \\
\hline
30 \textbf{*} & Rayakan completion dengan date night dan reflection bersama. \\
\hline
\end{tabular}

\vspace{0.5cm}
\textit{* Hari 30: Buat commitment untuk melanjutkan minimal 3 praktik yang paling berdampak.}
\end{exercisebox}

\section{Ringkasan dan Action Steps}

\begin{insightbox}
\textbf{Key Takeaways:}
\begin{enumerate}
    \item \highlight{Fondness adalah fondasi} --- tanpanya, konflik akan menghancurkan hubungan.
    \item Fondness bisa \highlight{dibangun kembali} meski hilang selama bertahun-tahun.
    \item Konsistensi lebih penting daripada grand gestures.
    \item Setiap attachment style memiliki bahasa fondness yang berbeda.
    \item Fondness harus diungkapkan \textbf{setiap hari}, bukan hanya saat senang.
\end{enumerate}
\end{insightbox}

\subsection{Action Plan untuk 7 Hari ke Depan}

\begin{exercisebox}[Langkah Langsung]
\textbf{Hari Ini:}
\begin{itemize}
    \item Kerjakan Fondness Assessment Scale
    \item Diskusikan hasil dengan pasangan
\end{itemize}

\textbf{Besok:}
\begin{itemize}
    \item Mulai dengan 1 apresiasi verbal spesifik
    \item Ciptakan 1 micro-moment affection
\end{itemize}

\textbf{Minggu Ini:}
\begin{itemize}
    \item Pilih 3 praktik dari 20+ tindakan yang akan dilakukan setiap hari
    \item Set reminder untuk consistency
    \item Mulai 30-Day Challenge
\end{itemize}

\textbf{Minggu Depan:}
\begin{itemize}
    \item Review: Apa yang berhasil? Apa yang perlu disesuaikan?
    \item Naikkan tingkat kesulitan/depth dari ungkapan
    \item Jadwalkan quality time khusus
\end{itemize}
\end{exercisebox}

\begin{tipbox}[Mantra Harian]
\begin{center}
\Large
``\textit{Hari ini, aku akan mencari dan mengungkapkan\\
setidaknya satu hal yang aku hargai dari pasanganku.}''
\end{center}
\end{tipbox}
